\documentclass[12pt,a4paper]{report}

\usepackage[utf8x]{inputenc}
\usepackage{amsmath}
\usepackage{amsfonts}
\usepackage{amssymb}
%\usepackage[linesnumbered,lined,boxed,commentsnumbered]{algorithm2e}
\usepackage{tikz}
\usetikzlibrary{arrows}

\begin{document}

\begin{titlepage}
	\centering
	{\scshape\LARGE Universidad Nacional Autónoma de México \par}
	\vspace{1cm}
	{\scshape\Large Computación Distribuida\par}
	\vspace{1.5cm}
	{\huge\bfseries Tarea 3\par}
	\vspace{.5cm}
	{\Large\itshape Edgar Quiroz Castañeda \par}
    \vspace{.5cm}
	{\Large\itshape Jerónimo Almeida Rodríguez \par}
	\vfill
	 \includegraphics[width=0.5\textwidth]{escudo_f-ciencias.png}
	\vfill

% Bottom of the page
	{\large Jueves 30 de Agosto del 2018 \par}
\end{titlepage}

\pagebreak
\setlength{\voffset}{-0.75in}
\setlength{\headsep}{5pt}


\begin{enumerate}
	%Ejercicio 1
	\item {
	
		\item[c] {
			% Declare layers
			\pgfdeclarelayer{background}
			\pgfsetlayers{background,main}
			% Styles
			\tikzstyle{information text}=[text badly centered,font=\small,text 
											width=3cm]
			\begin{tikzpicture}[scale=.8,cap=round]
			% The graphic
			    \begin{scope}[>=stealth', line width=1pt]
			        \draw[->] (1,.9) node[below, information text]
            			{Phonological representation of root form } -- (1,1.8);
			        \draw[->] (5,-.2) node[below,information text]
            			{Wickelfeature representation of root form } -- (5,.8);
			        \draw[->] (11,-.2) node[below,information text]
            			{Wickelfeature representation of past tense } -- (11,.8);
			        \draw[->] (16,0.9) node[below,information text]
            			{Phonological representation of past tense } -- (16,1.8);
			    \end{scope}
			    \draw (3,6) node[information text] { Fixed Encoding Network };
			    \draw (8,6) node[information text, text width=4cm, ]
			        { Pattern Associator Modifiable Connections };
			    \draw (13.5,6) node[information text] { Decoding/Binding Network };
			    % draw the nodes
			    \foreach \x in {1,16}
			        \foreach \y in {2,3,4} {
				        \filldraw[fill=white] (\x,\y) circle (0.1);
			        }
				    \foreach \x in {5,11}
				        \foreach \y in {1,2,3,4,5} {
				            \filldraw[fill=white] (\x,\y) circle (0.1);
				        }
			    % The lines connecting the nodes are drawn in the background layer.
			    % This way we can hide the lines behind the nodes and don't worry
			    % about the width of each node.    
			    \begin{pgfonlayer}{background}
			        % we add the lines for the nodes starting in y 2,3, and 4
			        \foreach \xa / \xb in {1 / 5, 5 / 11 , 11 / 5 , 16 / 11}
			            \foreach \ya / \yb / \yc / \yd / \ye in {2 / 3 / 4 / 5 / 1, 
			            3 / 4 / 5 / 1 / 2, 4 / 5 / 1 / 2 / 3} {
			                \draw[->] (\xa,\ya) -- (\xb,\ya);
			                \draw (\xa,\ya) -- (\xb,\yb);
            			    \draw (\xa,\ya) -- (\xb,\yc);
			                \draw (\xa,\ya) -- (\xb,\yd);
            			    \draw (\xa,\ya) -- (\xb,\ye);
			            }
			        % add remaining lines from y1 to y5
			        \foreach \xa / \xb in {5 / 11 , 11 / 5}
			            \foreach \ya / \yb in {1 / 5, 5 / 1} {
            			\draw (\xa,\ya) -- (\xb,\ya);
			            \draw (\xa,\ya) -- (\xb,\yb);
			        }
			    \end{pgfonlayer}
			\end{tikzpicture}
		}
	
	}

	%Ejercicio 2
	\item {
		
	}
	


\end{enumerate}
\end{document}